\documentclass[11pt]{article}
\usepackage[utf8]{inputenc}
\usepackage{amsmath}
\usepackage{amsfonts}
\usepackage{amssymb}
\usepackage{graphicx}
\usepackage[super]{nth}
\usepackage{amsthm}
\usepackage{bm}
\newtheorem{theorem}{Theorem}
\newtheorem{objective}{Objective}
\newtheorem{model}{Model}
\usepackage{xcolor}
\definecolor{light-gray}{gray}{0.95}
\newcommand{\code}[1]{\colorbox{light-gray}{\texttt{#1}}}
\usepackage{listings}
\usepackage{placeins}
\usepackage{enumitem}
\usepackage[T1]{fontenc}

\usepackage[authoryear]{natbib}

\lstset{language=R}


\makeatletter
\renewcommand{\maketitle}{
\begin{center}

\pagestyle{empty}

{\Large \bf \@title\par}
\vspace{1cm}

{\LARGE Marcel Gietzmann-Sanders}\\[1cm]

STAT651 - Statistical Theory I \\
University of Alaska Fairbanks


\end{center}
}\makeatother


\title{Homework 8}

\date{2025}
\setcounter{tocdepth}{2}
\begin{document}
\maketitle

\section*{Question 1}

\subsection*{a)}

Assuming we know the number of red $X$ and blue $Y$ balls we got, then the number of sets of balls that satisfy this combination of red and blue is:

$$N(X, Y) = {7 \choose X}{2 \choose Y}{4 \choose 3 - X - Y}$$

This is because our three tasks would be choose $X$ balls from a set of 7, choose $Y$ balls from a set of $2$, and then choose the remaining number of balls ($3-X-Y$) from 4.

\begin{center}
\begin{tabular}{ c| c c c c }
$N(X,Y)$ & $X=0$ & $X=1$ & $X=2$ & $X=3$ \\
 \hline
$Y=0$ & 4 & 42 & 84 & 35 \\ 
$Y=1$ & 12 & 56 & 42 &  \\
$Y=2$ & 4 & 7 &  & 
\end{tabular}
\end{center}

For a total of 286 different sets of three balls. Therefore our PMF is:

\begin{center}
\begin{tabular}{ c| c c c c }
$f(X,Y)$ & $X=0$ & $X=1$ & $X=2$ & $X=3$ \\
 \hline
$Y=0$ & 4/286 & 42/286 & 84/286 & 35/286 \\ 
$Y=1$ & 12/286 & 56/286 & 42/286 &  \\
$Y=2$ & 4/286 & 7/286 &  & 
\end{tabular}
\end{center}

\subsection*{b) and c)}

\begin{center}
\begin{tabular}{ c| c c c c | c}
$f(X,Y)$ & $X=0$ & $X=1$ & $X=2$ & $X=3$ & $f(Y)$ \\
 \hline
$Y=0$ & 4/286 & 42/286 & 84/286 & 35/286 & 165/286\\ 
$Y=1$ & 12/286 & 56/286 & 42/286 &  & 110/286 \\
$Y=2$ & 4/286 & 7/286 &  &  & 11/286 \\
\hline
$f(X)$ & 20/286 & 105/286 & 126/286 & 35/286 & 
\end{tabular}
\end{center}

\section*{Question 2}

\subsection*{a)}

The marginal distribution of $X$ is given in the question as is just:

\begin{center}
\begin{tabular}{c| c}
$P(X=0)$ & 0.4 \\
$P(X=1)$ & 0.3 \\
$P(X=2)$ & 0.2 \\
$P(X=3)$ & 0.1
\end{tabular}
\end{center}

\subsection*{b)}

In this case there are two people in line asking a total of 5 questions. This means one of them is asking 2 questions and the other is asking three. Let us call the event that the first person asked 2 questions and the second 3 $A$ and the event that the first person asked 3 and the second 2 $B$. Then the event $C$ that 5 questions were asked is just $A\cup B$. So $P(A\cup B) = P(A) + P(B) - P(A\cap B)=P(A) + P(B)$ given the intersection between these two events does not exist. $P(A) + P(B) = 0.5(0.1) + 0.1(0.5) = 0.1$. Therefore $P(Y=5 | X=2)=0.1$. So $f(2,5)=P(Y=5 | X=2)P(X=2)=0.1(0.2)=0.02$. 

\subsection*{c)}

In this case we know the number of people in line is 3. So we want to find $P(Y=6 | X=3)$. Our cases are:

\begin{center}
\begin{tabular}{c c c | c}
first & second & third & p \\
\hline
1 & 2 & 3 & $$0.4(0.5)0.1=0.02$$\\ 
1 & 3 & 2 & 0.02 \\ 
2 & 1 & 3 & 0.02 \\ 
2 & 3 & 1 & 0.02 \\ 
3 & 1 & 2 & 0.02 \\ 
3 & 2 & 1 & 0.02

\end{tabular}
\end{center}

As there are no intersections between these events the cumulative probability across their union is just their sum which is $0.12$. Therefore $P(Y=6 | X=3)=0.12$ and so $f(3,6)=P(Y=6 | X=3)P(X=3)=0.12(0.1)=0.012$. 

\section*{Question 3}

\subsection*{a)}

\subsection*{b)}

Given this is an exponent we know it'll always be positive. So question is does this thing integrate to one?

$$\int_0^{\infty}\int_0^{y}2e^{-(x+y)} dxdy=-2\int_0^{\infty}e^{-(x+y)} \Big|_0^y dy$$
$$=2\int_0^{\infty}e^{-y}-e^{-2y} dy=2\left(e^{-2y}/2-e^{-y}) \right)\Big|_0^{\infty}=1$$

So it is indeed a pdf!

\section*{Question 4}
\subsection*{a)}

\subsection*{b)}

We are taking the sum of two positive numbers and will therefore get a positive number. Time to find out if the function integrates to one:

$$\int_0^1 \int_0^x2(x+y)dydx=2\int_0^1xy+y^2/2\Big|_0^x dx$$
$$=2\int_0^1 x^2dx=1$$










\end{document}